\documentclass[../main.tex]{subfiles}
\begin{document}
%=========================================================
% TĐ_PBT_IQB_HK1_Tuan06_OnTapPhepNhan_ToanLoiVan
% Phiếu bài tập Toán tư duy cấp độ B
%   Buổi 1: ÔN TẬP PHÉP NHÂN
%   Buổi 2: TOÁN CÓ LỜI VĂN -- PHÉP CỘNG, PHÉP TRỪ
% Đi kèm giáo án GiaoAn_Tuan06_OnTapPhepNhan_ToanLoiVan (số bài khớp với giáo án)
% Soạn theo mẫu (\hktuan, bt, listEX)
% Hình đặt trong thư mục Hinh_Tuan06_PBT/ (cùng cấp với file này), tên hình có tiền tố t06_
% Đường dẫn tính từ thư mục của main.tex khi biên dịch main.tex
\graphicspath{{Hinh_GiaoAn_Tuan06/}{Data_PhieuBaiTap_HK0/Hinh_GiaoAn_Tuan06/}{Hinh_Tuan06_PBT/}{Data_PhieuBaiTap_HK0/Hinh_Tuan06_PBT/}}
%=========================================================

%---------- Màu dùng riêng cho phiếu ----------
\definecolor{TTcam}{HTML}{F39C34}
\definecolor{TTcamdam}{HTML}{C9731A}
\definecolor{TTxanh}{HTML}{4F86C6}
\definecolor{TTxanhdam}{HTML}{2D5F99}
\definecolor{TTdo}{HTML}{C4313A}

%---------- Lệnh dùng riêng cho phiếu ----------
% Chèn hình căn giữa, dính liền với đề bài phía trên (không bị tách trang)
% Tham số tùy chọn = chiều cao tối đa của hình (mặc định 8cm); hình không rộng quá khổ giấy
\providecommand{\TThinh}{}\renewcommand{\TThinh}[2][8cm]{\par\nopagebreak\vspace{2pt}\nopagebreak{\centering\includegraphics[height=#1,width=\linewidth,keepaspectratio]{#2}\par}\vspace{2pt}}
% Chỗ chấm điền
\providecommand{\TTcham}[1][3cm]{\makebox[#1]{\dotfill}}
% Ô vuông trống để điền dấu hoặc Đ/S
\providecommand{\TTo}{}\renewcommand{\TTo}{\hspace{3pt}\framebox[0.75cm]{\rule{0pt}{0.42cm}}\hspace{3pt}}
% Dòng thông tin học sinh
\providecommand{\TTthongtin}{%
  \noindent\begin{tabular}{|p{7.6cm}|p{4cm}|p{4cm}|}
    \hline
    \textbf{Họ và Tên HS:}\dotfill & \textbf{Lớp:}\dotfill & \textbf{Ngày:}\dotfill \\
    \hline
  \end{tabular}\par\medskip}
% Đánh lại số bài từ 1 cho buổi 2 (chỉ khi môi trường bt có bộ đếm tên bt)
\makeatletter
\providecommand{\TTdatlaibai}{\@ifundefined{c@bt}{}{\setcounter{bt}{0}}}
\makeatother
% Khung "Bài giải" có n dòng kẻ chấm: \TTbaigiai{n}
\providecommand{\TTbaigiai}{}\renewcommand{\TTbaigiai}[1]{\par\nopagebreak\vspace{4pt}{\centering\underline{\textbf{Bài giải}}\par}\nopagebreak
  \foreach \i in {1,...,#1}{\vspace{9pt}\noindent\makebox[\linewidth]{\dotfill}\par}\medskip}
% n dòng kẻ chấm (không có chữ "Bài giải"): \TTdongcham{n}
\providecommand{\TTdongcham}{}\renewcommand{\TTdongcham}[1]{\par\nopagebreak
  \foreach \i in {1,...,#1}{\vspace{9pt}\noindent\makebox[\linewidth]{\dotfill}\par}\medskip}
% Ô li nhỏ để đặt tính (4 cột x 4 hàng)
\providecommand{\TTluoi}{}\renewcommand{\TTluoi}{\tikz\draw[black!25,step=0.55cm] (0,0) grid (2.2,2.2);}
% Một ô hình ở Bài 1 (Buổi 1): \TTohinh{nhãn}{tên hình}{chiều cao}
\providecommand{\TTohinh}{}\renewcommand{\TTohinh}[3]{%
  \begin{minipage}[t]{0.48\linewidth}
    #1\par\nopagebreak
    \parbox[c][3.1cm][c]{\linewidth}{\centering\includegraphics[height=#3,width=\linewidth,keepaspectratio]{#2}}\par
    Phép cộng: \TTcham[5.6cm]\par\vspace{6pt}
    Phép nhân: \TTcham[5.6cm]
  \end{minipage}}
% Thẻ số (TikZ)
\tikzset{TTthe/.style={draw=black!30,fill=white,rounded corners=4pt,minimum width=1.2cm,minimum height=1.45cm,font=\huge\bfseries,text=mauchinh}}
\renewcommand{\lop}{B}
%%%%%%%%%%%%%%%%%%%%%%%%%%%%%%%%%%%%%%%%%%%%%%%%%%%%%%%%%%
%                       BUỔI 1
%%%%%%%%%%%%%%%%%%%%%%%%%%%%%%%%%%%%%%%%%%%%%%%%%%%%%%%%%%
%\setcounter{section}{3} % không chỉnh sửa
\hktuan{1}{14}
\par\medskip % xuống dòng: tiêu đề không nằm chung dòng với khung HK1 - Tuần 6
{\centering
        {\color{mauchinh!120}\bfseries\Large PHIẾU BÀI TẬP TOÁN TƯ DUY CẤP ĐỘ B}\\[3pt]
        {\color{mauchinh!120}\bfseries\Large CHỦ ĐỀ: ÔN TẬP PHÉP NHÂN (BUỔI 1)}\par}
\vspace{6pt}
\TTthongtin

\begin{bt}
    Quan sát hình rồi viết phép cộng và phép nhân thích hợp.
    \par\nopagebreak\vspace{4pt}\nopagebreak
        \noindent\TTohinh{a) Có 3 lốc sữa, mỗi lốc 4 hộp.}{ga06_locsua.png}{1.6cm}\hfill
    \TTohinh{b) Có 4 lọ hoa, mỗi lọ 5 bông.}{ga06_lohoa.png}{2.9cm}\par\vspace{12pt}
    \noindent\TTohinh{c) Có 5 đĩa cam, mỗi đĩa 3 quả.}{ga06_diacam.png}{1.6cm}\hfill
    \TTohinh{d) Có 4 hộp bút, mỗi hộp 8 chiếc.}{ga06_hopbut.png}{2.7cm}\par
\end{bt}

\begin{bt}
    Nối mỗi tổng với phép nhân thích hợp.
    \par\nopagebreak\vspace{6pt}\nopagebreak{\centering
    \begin{tikzpicture}[
        tong/.style={draw=TTcamdam,fill=TTcam!12,thick,rounded corners=6pt,minimum width=2.9cm,minimum height=0.95cm,font=\large},
        nhan/.style={draw=TTxanhdam,fill=TTxanh!10,thick,rounded corners=6pt,minimum width=2.9cm,minimum height=0.95cm,font=\large}]
        \foreach \x/\t [count=\k] in {0/{2 + 2 + 2},3.4/{5 + 5 + 5 + 5},6.8/{3 + 3},10.2/{6 + 6 + 6},13.6/{4 + 4 + 4 + 4}}{
            \node[tong] (t\k) at (\x,2.6) {$\t$};
            \fill (t\k.south) ++(0,-0.12) circle (2.4pt);}
        \foreach \x/\t [count=\k] in {0/{5 \times 4},3.4/{6 \times 3},6.8/{2 \times 3},10.2/{4 \times 4},13.6/{3 \times 2}}{
            \node[nhan] (n\k) at (\x,0) {$\t$};
            \fill (n\k.north) ++(0,0.12) circle (2.4pt);}
    \end{tikzpicture}\par}
\end{bt}

\begin{bt}
    Viết số thích hợp vào chỗ trống (theo mẫu).
    \par\nopagebreak\vspace{4pt}\nopagebreak{\centering
        \renewcommand{\arraystretch}{1.6}
        \begin{tabular}{|>{\centering\arraybackslash}p{4cm}|>{\centering\arraybackslash}p{2.8cm}|>{\centering\arraybackslash}p{2.8cm}|>{\centering\arraybackslash}p{2.8cm}|}
            \hline
            \textbf{Phép nhân} & \textbf{Thừa số} & \textbf{Thừa số} & \textbf{Tích} \\ \hline
            $6 \times 4 = 24$ & \textit{6} & \textit{4} & \textit{24} \\ \hline
            $7 \times 3 = 21$ & & & \\ \hline
            $9 \times 5 = 45$ & & & \\ \hline
            \TTcham[3cm] & 8 & 2 & \\ \hline
            \TTcham[3cm] & 5 & 6 & \\ \hline
        \end{tabular}\par}
\end{bt}
\newpage
\begin{bt}
    Tính nhẩm.
    \begin{multicols}{4}
    \begin{listEX}
        \item $2 \times 7 =$ \TTcham[1.1cm]
        \item $3 \times 6 =$ \TTcham[1.1cm]
        \item $4 \times 8 =$ \TTcham[1.1cm]
        \item $5 \times 9 =$ \TTcham[1.1cm]
        \item $6 \times 7 =$ \TTcham[1.1cm]
        \item $7 \times 8 =$ \TTcham[1.1cm]
        \item $8 \times 6 =$ \TTcham[1.1cm]
        \item $9 \times 9 =$ \TTcham[1.1cm]
        \item $6 \times 1 =$ \TTcham[1.1cm]
        \item $1 \times 9 =$ \TTcham[1.1cm]
        \item $7 \times 0 =$ \TTcham[1.1cm]
        \item $0 \times 5 =$ \TTcham[1.1cm]
        \item $8 \times 9 =$ \TTcham[1.1cm]
        \item $9 \times 4 =$ \TTcham[1.1cm]
        \item $5 \times 5 =$ \TTcham[1.1cm]
        \item $7 \times 7 =$ \TTcham[1.1cm]
    \end{listEX}
    \end{multicols}
\end{bt}

\begin{bt}
    Điền số thích hợp vào chỗ chấm.
    \begin{multicols}{4}
    \begin{listEX}
        \item $6 \times$ \TTcham[0.9cm] $= 30$
        \item \TTcham[0.9cm] $\times\, 4 = 28$
        \item $9 \times$ \TTcham[0.9cm] $= 9$
        \item \TTcham[0.9cm] $\times\, 8 = 0$
    \end{listEX}
    \end{multicols}
\end{bt}
\begin{bt}
    Viết tiếp số thích hợp vào chỗ chấm theo quy luật.
    \begin{listEX}
        \item $2,\ 4,\ 6,\ 8,$ \TTcham[1cm], \TTcham[1cm], $14,$ \TTcham[1cm]
        \item $5,\ 10,\ 15,$ \TTcham[1cm], $25,$ \TTcham[1cm], \TTcham[1cm], $40$
        \item $3,\ 6,\ 9,$ \TTcham[1cm], \TTcham[1cm], $18,$ \TTcham[1cm], $24$
        \item $36,\ 32,\ 28,$ \TTcham[1cm], \TTcham[1cm], $16,$ \TTcham[1cm]
    \end{listEX}
\end{bt}
\begin{bt}
    Điền dấu $>$, $<$, $=$ vào ô trống mà không cần tính. Giải thích cho bạn bên cạnh vì sao con điền dấu đó.
    \begin{multicols}{2}
    \begin{listEX}
        \item $4 \times 5$ \TTo\ $5 \times 4$
        \item $7 \times 6$ \TTo\ $7 \times 8$
        \item $13 \times 3$ \TTo\ $12 \times 3$
        \item $5 + 5 + 5 + 5$ \TTo\ $5 \times 3$
        \item $9 \times 0$ \TTo\ $1 \times 0$
        \item $8 \times 1$ \TTo\ $8 + 1$
        \item $9 \times 6$ \TTo\ $6 \times 9$
        \item $25 \times 2$ \TTo\ $25 \times 4$
    \end{listEX}
    \end{multicols}
\end{bt}

\begin{bt}
    Trong mỗi cặp, khoanh vào phép nhân có tích lớn hơn (không cần tính).
    \begin{multicols}{2}
    \begin{listEX}
        \item \makebox[2.4cm][l]{$6 \times 7$}\makebox[2.4cm][l]{$6 \times 9$}
        \item \makebox[2.4cm][l]{$25 \times 4$}\makebox[2.4cm][l]{$24 \times 4$}
        \item \makebox[2.4cm][l]{$18 \times 3$}\makebox[2.4cm][l]{$18 \times 2$}
        \item \makebox[2.4cm][l]{$7 \times 5$}\makebox[2.4cm][l]{$5 \times 8$}
        \item \makebox[2.4cm][l]{$4 + 4 + 4$}\makebox[2.4cm][l]{$4 \times 4$}
        \item \makebox[2.4cm][l]{$9 \times 0$}\makebox[2.4cm][l]{$1 \times 9$}
    \end{listEX}
    \end{multicols}
\end{bt}

\begin{bt}
    Trả lời các câu hỏi sau.
    \begin{listEX}
        \item Một phép nhân có thừa số thứ nhất là 6, thừa số thứ hai là 5. Nếu tăng thừa số thứ hai thêm 2 đơn vị thì tích tăng thêm bao nhiêu đơn vị?
        \TTdongcham{2}
        \item Cho phép nhân $8 \times 5 = 40$. Nếu thay số 5 bằng số 3 thì tích tăng hay giảm bao nhiêu đơn vị?
        \TTdongcham{2}
        \item Trong một phép nhân, nếu tăng thừa số thứ nhất thêm 3 đơn vị thì tích tăng 15 đơn vị. Tìm thừa số thứ hai.
        \TTdongcham{2}
    \end{listEX}
\end{bt}

\begin{bt}
    Đúng ghi Đ, sai ghi S vào ô trống.
    \par\nopagebreak\vspace{4pt}\nopagebreak{\centering
        \renewcommand{\arraystretch}{1.55}
        \begin{tabular}{|>{\raggedright\arraybackslash}p{12.4cm}|>{\centering\arraybackslash}p{1.4cm}|}
            \hline
            a) $7 \times 4 = 4 \times 7$ & \\ \hline
            b) $9 \times 0 = 9$ & \\ \hline
            c) $1 \times 6 = 6$ & \\ \hline
            d) $8 \times 5 < 8 \times 4$ & \\ \hline
            e) $3 \times 5 = 3 + 3 + 3 + 3 + 3$ & \\ \hline
            f) Trong phép nhân $6 \times 3 = 18$, số 18 được gọi là thừa số. & \\ \hline
        \end{tabular}\par}
\end{bt}

\begin{bt}
    Mỗi chiếc ô tô con có 4 bánh xe. Hỏi 6 chiếc ô tô như thế có bao nhiêu bánh xe?
    \TTbaigiai{3}
\end{bt}

\begin{bt}
    Mỗi tuần lễ có 7 ngày. Hỏi 3 tuần lễ có bao nhiêu ngày?
    \TTbaigiai{3}
\end{bt}

\begin{bt}
    Mỗi đĩa có 3 quả cam. Ngày Tết, bà bày 9 đĩa như thế. Hỏi bà bày tất cả bao nhiêu quả cam?
    \TTbaigiai{3}
\end{bt}

\begin{bt}
    Mỗi chuồng có 5 con thỏ. Hỏi 7 chuồng như thế có tất cả bao nhiêu con thỏ?
    \TTbaigiai{3}
\end{bt}

\begin{bt}
    Cho phép nhân $5 \times 8 = 40$. Nếu đổi số 8 thành số 6 thì tích tăng hay giảm bao nhiêu đơn vị?
    \TTdongcham{3}
\end{bt}

\begin{bt}
    Mỗi hộp có 6 cái bánh. Mẹ mua 4 hộp như thế. Cả nhà đã ăn hết 9 cái bánh. Hỏi còn lại bao nhiêu cái bánh?
    \TTbaigiai{5}
\end{bt}

\begin{bt}
    Trong sân có 5 chiếc xe đạp và 4 chiếc xe ba bánh. Mỗi xe đạp có 2 bánh, mỗi xe ba bánh có 3 bánh. Hỏi trong sân có tất cả bao nhiêu bánh xe?
    \TTbaigiai{6}
\end{bt}
\newpage
\begin{bt}
    Mỗi bàn học có 2 bạn ngồi. Lớp 2A có 9 bàn. Cô giáo kê thêm 3 bàn như thế.
    \begin{listEX}
        \item Lớp có thể xếp thêm bao nhiêu bạn ngồi?
        \item Bây giờ lớp có thể xếp tất cả bao nhiêu bạn ngồi?
    \end{listEX}
    \TTbaigiai{7}
\end{bt}

\begin{bt}
    Nhà Lan nuôi 4 chuồng thỏ, mỗi chuồng 6 con. Nhà Hoa nuôi 3 chuồng thỏ, mỗi chuồng 8 con. Nhà nào nuôi nhiều thỏ hơn? Nhiều hơn mấy con?
    \TTbaigiai{7}
\end{bt}
%---------- PHÁT TRIỂN TƯ DUY (Buổi 1) ----------
% Dán đoạn này ngay SAU phần "BÀI TOÁN NÂNG CAO" của Buổi 1, TRƯỚC khối "%  BUỔI 2"
% Ghi chú nguồn dạng toán (chữ nghiêng, màu xám)
\providecommand{\TTnguon}{}\renewcommand{\TTnguon}[1]{{\small\itshape\color{black!55}(#1)}}

\begin{bt}
    Mỗi con bọ rùa có 6 chân. Trên một chiếc lá có một số con bọ rùa. Bạn Na đếm được tất cả 18 cái chân. Hỏi trên chiếc lá có mấy con bọ rùa?
    \TTbaigiai{3}
\end{bt}

\begin{bt}
    Mỗi hình giống nhau thay cho cùng một số. Em hãy tìm số thích hợp.
    \par\nopagebreak\vspace{6pt}\nopagebreak{\centering\Large
        $\bigstar + \bigstar + \bigstar = 12$ \qquad\qquad $\bullet = \bigstar \times 2$\par}
    \vspace{8pt}
    \begin{listEX}
        \item $\bigstar =$ \TTcham[1.5cm]
        \item $\bullet =$ \TTcham[1.5cm]
        \item $\bigstar \times \bullet =$ \TTcham[1.5cm]
    \end{listEX}
\end{bt}


%%%%%%%%%%%%%%%%%%%%%%%%%%%%%%%%%%%%%%%%%%%%%%%%%%%%%%%%%%
%                       BUỔI 2
%%%%%%%%%%%%%%%%%%%%%%%%%%%%%%%%%%%%%%%%%%%%%%%%%%%%%%%%%%
\newpage
\TTdatlaibai
{\centering
        {\color{mauchinh!120}\bfseries\Large PHIẾU BÀI TẬP TOÁN TƯ DUY CẤP ĐỘ B}\\[3pt]
        {\color{mauchinh!120}\bfseries\Large CHỦ ĐỀ: TOÁN CÓ LỜI VĂN -- PHÉP CỘNG,\\[2pt] PHÉP TRỪ (BUỔI 2)}\par}
\vspace{6pt}
\TTthongtin
%---------- Lệnh cho Bài 1, Bài 2 (Buổi 2) ----------
% Bài 1: số -- ô trống -- số -- dấu = -- chỗ chấm, mỗi phần rộng cố định nên thẳng cột
\providecommand{\TTdien}{}\renewcommand{\TTdien}[2]{\hfill\makebox[0.9cm][r]{$#1$}\TTo\makebox[0.75cm][l]{$#2$}\makebox[0.45cm]{$=$}\TTcham[1.6cm]}
% Bài 2: ô li tiểu học để đặt tính: 4 x 4 ô vuông lớn 8 mm (nét đậm),
% mỗi ô lớn chia 4 x 4 ô li nhỏ 2 mm (nét mảnh)
\providecommand{\TTluoilon}{}\renewcommand{\TTluoilon}{\tikz{%
  \draw[black!20,line width=0.3pt,step=2mm] (0,0) grid (3.2,3.2);%
  \draw[black!55,line width=0.5pt,step=8mm] (0,0) grid (3.2,3.2);%
  \draw[black!55,line width=0.5pt] (0,0) rectangle (3.2,3.2);}}

\begin{bt}
    Viết dấu $+$ hoặc $-$ vào ô trống rồi tính.
    \begin{listEX}
        \item Trên cây có 9 quả na, bà hái xuống 4 quả. \TTdien{9}{4}
        \item Hộp có 6 bút chì, bố mua thêm 5 bút chì. \TTdien{6}{5}
        \item Lớp có 14 bạn nam và 13 bạn nữ. \TTdien{14}{13}
        \item Có 20 quả bóng bay, 8 quả bị vỡ. \TTdien{20}{8}
        \item Đĩa có 15 cái bánh, em ăn 3 cái. \TTdien{15}{3}
        \item Ao có 12 con vịt, có thêm 6 con xuống ao. \TTdien{12}{6}
    \end{listEX}
\end{bt}

\begin{bt}
    Đặt tính rồi tính.
    \par\nopagebreak\vspace{4pt}\nopagebreak{\centering
        \setlength{\tabcolsep}{0.7cm}
        \begin{tabular}{ccc}
            $34 + 25$ & $48 + 17$ & $25 + 8$ \\[3pt]
            \TTluoilon & \TTluoilon & \TTluoilon \\[12pt]
            $63 - 27$ & $45 - 12$ & $90 - 36$ \\[3pt]
            \TTluoilon & \TTluoilon & \TTluoilon \\ [12pt]
             $45 - 28$ & $81 - 12$ & $73 + 17$ \\[3pt]
            \TTluoilon & \TTluoilon & \TTluoilon \\
        \end{tabular}\par}
\end{bt}
\newpage
\begin{bt}
    Tủ sách của lớp có 48 quyển truyện. Các bạn góp thêm 17 quyển. Hỏi tủ sách của lớp có tất cả bao nhiêu quyển truyện?
    \TTbaigiai{3}
\end{bt}

\begin{bt}
    Buổi sáng, cửa hàng bán được 34 quả bưởi. Buổi chiều, cửa hàng bán được 25 quả bưởi. Hỏi cả ngày cửa hàng bán được bao nhiêu quả bưởi?
    \TThinh[3.2cm]{ga06_robuoi.png}
    \par\nopagebreak{\centering\begin{tabular}{p{0.45\linewidth}p{0.45\linewidth}}\centering Buổi sáng: \textbf{34 quả} & \centering\arraybackslash Buổi chiều: \textbf{25 quả}\end{tabular}\par}
    \noindent\begin{minipage}[t]{0.36\linewidth}
        \textbf{Tóm tắt}
        \TTdongcham{3}
    \end{minipage}\hfill
    \begin{minipage}[t]{0.6\linewidth}
        \TTbaigiai{3}
    \end{minipage}
\end{bt}

\begin{bt}
    Trên dây điện có 12 con chim sẻ đang đậu. Có 5 con bay đi. Hỏi trên dây điện còn lại bao nhiêu con chim sẻ?
    \TThinh[3cm]{ga06_chimse.png}
    \TTbaigiai{3}
\end{bt}

\begin{bt}
    Cửa hàng có 63 kg gạo, đã bán được 27 kg. Hỏi cửa hàng còn lại bao nhiêu ki-lô-gam gạo?
    \TTbaigiai{3}
\end{bt}
\begin{bt}
    Nhà bác Ba nuôi 42 con gà. Bác bán đi 15 con. Hỏi nhà bác Ba còn lại bao nhiêu con gà?
    \TTbaigiai{3}
\end{bt}

\begin{bt}
    Tổ Một trồng được 26 cây. Tổ Hai trồng được nhiều hơn tổ Một 7 cây. Hỏi tổ Hai trồng được bao nhiêu cây?\par
    Viết số thích hợp vào chỗ chấm trong phần tóm tắt rồi giải bài toán.
    \par\nopagebreak\vspace{4pt}\nopagebreak
    \noindent\begin{minipage}[t]{0.42\linewidth}
        \textbf{Tóm tắt}\par\vspace{9pt}
        Tổ Một: \TTcham[1.2cm] cây\par\vspace{9pt}
        Tổ Hai nhiều hơn tổ Một: \TTcham[1.2cm] cây\par\vspace{9pt}
        Tổ Hai: \TTcham[1.2cm] cây?
    \end{minipage}\hfill
    \begin{minipage}[t]{0.54\linewidth}
        \TTbaigiai{3}
    \end{minipage}
\end{bt}

\begin{bt}
    Anh cao 95 cm. Em thấp hơn anh 18 cm. Hỏi em cao bao nhiêu xăng-ti-mét?\par
    Viết tóm tắt rồi giải bài toán.
    \par\nopagebreak\vspace{4pt}\nopagebreak
    \noindent\begin{minipage}[t]{0.42\linewidth}
        \textbf{Tóm tắt}
        \TTdongcham{3}
    \end{minipage}\hfill
    \begin{minipage}[t]{0.54\linewidth}
        \TTbaigiai{3}
    \end{minipage}
\end{bt}
\newpage
\begin{bt}
    Nhìn tóm tắt, em hãy nêu bài toán rồi giải bài toán đó.
    \par\nopagebreak\vspace{4pt}\nopagebreak
    \noindent\begin{minipage}[t]{0.42\linewidth}
        \textbf{Tóm tắt}\par\vspace{6pt}
        Thùng cam: 40 quả\par\vspace{6pt}
        Thùng quýt ít hơn thùng cam: 15 quả\par\vspace{6pt}
        Thùng quýt: \dots\ quả?
    \end{minipage}\hfill
    \begin{minipage}[t]{0.54\linewidth}
        \textbf{Bài toán}
        \TTdongcham{3}
    \end{minipage}
    \TTbaigiai{3}
\end{bt}
\begin{bt}
    Bao gạo thứ nhất nặng 35 kg. Bao gạo thứ hai nặng hơn bao thứ nhất 18 kg. Hỏi bao gạo thứ hai nặng bao nhiêu ki-lô-gam?
    \TTbaigiai{3}
\end{bt}

\begin{bt}
    Xe buýt có 28 hành khách. Đến bến, có 9 người xuống xe và 6 người lên xe. Hỏi lúc này trên xe có bao nhiêu hành khách?
    \TTbaigiai{5}
\end{bt}

\begin{bt}
    Tổ Một gấp được 34 ngôi sao. Tổ Hai gấp được ít hơn tổ Một 9 ngôi sao. Hỏi cả hai tổ gấp được bao nhiêu ngôi sao?
    \TTbaigiai{5}
\end{bt}

\begin{bt}
    Sau khi mẹ mua thêm 15 quả trứng thì trong tủ lạnh có tất cả 40 quả trứng. Hỏi lúc đầu trong tủ lạnh có bao nhiêu quả trứng?
    \TTbaigiai{3}
\end{bt}

\begin{bt}
    Bình cho bạn 12 cái kẹo thì Bình còn lại 18 cái kẹo. Hỏi lúc đầu Bình có bao nhiêu cái kẹo?
    \TTbaigiai{3}
\end{bt}

\begin{bt}
    Bao gạo nặng 45 kg. Bao gạo nặng hơn bao ngô 18 kg. Hỏi bao ngô nặng bao nhiêu ki-lô-gam?\par
    \TTbaigiai{5}
\end{bt}

\begin{bt}
    Cây bút chì dài 17 cm. Cây thước kẻ dài hơn bút chì 13 cm. Cục tẩy ngắn hơn bút chì 12 cm.
    \begin{listEX}
        \item Cây thước kẻ dài bao nhiêu xăng-ti-mét?
        \item Cục tẩy dài bao nhiêu xăng-ti-mét?
    \end{listEX}
    \TTbaigiai{5}
\end{bt}
%---------- PHÁT TRIỂN TƯ DUY (Buổi 2) ----------
% Dán đoạn này ngay SAU phần "BÀI TOÁN NÂNG CAO" của Buổi 2, TRƯỚC khối \begin{center} ... Hết ... \end{center}
% Ghi chú nguồn dạng toán (chữ nghiêng, màu xám)
\providecommand{\TTnguon}{}\renewcommand{\TTnguon}[1]{{\small\itshape\color{black!55}(#1)}}
\newpage
\begin{bt}
    Trên cành cây có một số con chim đậu. Có 6 con bay đi, sau đó có 4 con khác bay đến. Lúc này trên cành có 10 con chim. Hỏi lúc đầu trên cành có bao nhiêu con chim?
    \TTbaigiai{5}
\end{bt}

\begin{bt}
    Các bạn xếp thành một hàng dọc. Bạn An đứng thứ 5 tính từ đầu hàng và cũng đứng thứ 5 tính từ cuối hàng. Hỏi hàng đó có bao nhiêu bạn?
    \TTdongcham{5}
\end{bt}

\begin{bt}
    Năm nay mẹ 32 tuổi, con 7 tuổi.
    \begin{listEX}
        \item Năm nay mẹ hơn con bao nhiêu tuổi?
        \item Sau 3 năm nữa, mẹ hơn con bao nhiêu tuổi?
    \end{listEX}
    \TTbaigiai{7}
\end{bt}


\begin{center}
       \faGg\faGg\quad Hết \quad\faGg\faGg
\end{center}
\end{document}
